% GENERATED FILE. Edit canonical lessons, then run npm run generate:books.
% canonical-source-hash: fnv1a64:b9b87528d5e64895
% canonical-lessons: KA-C174-pravasi, KA-C174-sarovara, KA-C174-guhe, KA-C174-mrgalaya, KA-C174-aramane

\chapter{Tourist, Lake, Cave, Zoo, Palace}
\label{ch:ka-tourist-lake-cave-zoo-palace}
\hlchaptermodality{174}

\begin{chapteropening}
\textbf{By the end of this chapter:} \emph{I can say a tourist, a lake, a cave, a zoo and a palace.}

Complete the last lesson of chapter 174: I can say a tourist, a lake, a cave, a zoo and a palace.
\end{chapteropening}

\section[\texorpdfstring{\kn{ಪ್ರವಾಸಿ} (pravāsi)}{ಪ್ರವಾಸಿ (pravāsi)}]{\kn{ಪ್ರವಾಸಿ} (pravāsi) — a tourist}

\label{lesson:KA-C174-pravasi}



\begin{quote}
Before the new one: say the Kannada for a room, then the Kannada for a courtyard.
\end{quote}

\subsection*{You'll want to know: \kn{ಪ್ರವಾಸಿ}}
\textbf{\kn{ಪ್ರವಾಸಿ}} — \emph{pravāsi} — \textquotedblleft{}a tourist\textquotedblright{}.

\begin{cousinweb}
\textbf{\kn{ಪ್ರವಾಸಿ}} (\emph{pravāsi}) is Sanskrit \emph{pravāsin}, \textquotedblleft{}one who is away from home, a traveller\textquotedblright{}, from \emph{pravāsa}, \textquotedblleft{}a stay abroad\textquotedblright{}. Kannada drops the final \emph{-n}.
\end{cousinweb}

\begin{grammarlens}[title={What you've built}]
One more word for this chapter.
\end{grammarlens}

\subsection*{Your turn}
\begin{itemize}
\raggedright
  \item \emph{Say it:} \emph{pravāsi}
  \item \emph{Say it:} \emph{pravāsi}, once more
  \item \emph{Say it:} say \emph{pravāsi}
  \item \emph{Recall:} say the Kannada for a room, then the Kannada for a courtyard, then say \emph{pravāsi} again
\end{itemize}

\begin{tcolorbox}[breakable,colback=teal!4,colframe=teal!35!black,title={Before you move on}]
What does \kn{ಪ್ರವಾಸಿ} mean? (\textquotedblleft{}A tourist\textquotedblright{}.) Say it once more.
\end{tcolorbox}
\section[\texorpdfstring{\kn{ಸರೋವರ} (sarōvara)}{ಸರೋವರ (sarōvara)}]{\kn{ಸರೋವರ} (sarōvara) — a lake}

\label{lesson:KA-C174-sarovara}



\begin{quote}
Before the new one: say the Kannada for a courtyard, then the Kannada for a tourist.
\end{quote}

\subsection*{You'll want to know: \kn{ಸರೋವರ}}
\textbf{\kn{ಸರೋವರ}} — \emph{sarōvara} — \textquotedblleft{}a lake\textquotedblright{}.

\begin{cousinweb}
\textbf{\kn{ಸರೋವರ}} (\emph{sarōvara}) is Sanskrit \emph{sarovara}, \textquotedblleft{}a fine lake\textquotedblright{}: \emph{saras}, \textquotedblleft{}lake, pool\textquotedblright{} (CDIAL 13254), and \emph{vara}, \textquotedblleft{}best\textquotedblright{}.
\end{cousinweb}

\begin{grammarlens}[title={What you've built}]
One more word for this chapter.
\end{grammarlens}

\subsection*{Your turn}
\begin{itemize}
\raggedright
  \item \emph{Say it:} \emph{sarōvara}
  \item \emph{Say it:} \emph{sarōvara}, once more
  \item \emph{Say it:} say \emph{sarōvara}
  \item \emph{Recall:} say the Kannada for a courtyard, then the Kannada for a tourist, then say \emph{sarōvara} again
\end{itemize}

\begin{tcolorbox}[breakable,colback=teal!4,colframe=teal!35!black,title={Before you move on}]
What does \kn{ಸರೋವರ} mean? (\textquotedblleft{}A lake\textquotedblright{}.) Say it once more.
\end{tcolorbox}
\section[\texorpdfstring{\kn{ಗುಹೆ} (guhe)}{ಗುಹೆ (guhe)}]{\kn{ಗುಹೆ} (guhe) — a cave}

\label{lesson:KA-C174-guhe}



\begin{quote}
Before the new one: say the Kannada for a tourist, then the Kannada for a lake.
\end{quote}

\subsection*{You'll want to know: \kn{ಗುಹೆ}}
\textbf{\kn{ಗುಹೆ}} — \emph{guhe} — \textquotedblleft{}a cave\textquotedblright{}.

\begin{cousinweb}
\textbf{\kn{ಗುಹೆ}} (\emph{guhe}) is Sanskrit \emph{guhā}, \textquotedblleft{}hiding place, cave\textquotedblright{}, taken over whole as a learned loan, a tatsama, with Kannada's \emph{-e} for the Sanskrit \emph{-ā} (CDIAL 4220).
\end{cousinweb}

\begin{grammarlens}[title={What you've built}]
One more word for this chapter.
\end{grammarlens}

\subsection*{Your turn}
\begin{itemize}
\raggedright
  \item \emph{Say it:} \emph{guhe}
  \item \emph{Say it:} \emph{guhe}, once more
  \item \emph{Say it:} say \emph{guhe}
  \item \emph{Recall:} say the Kannada for a tourist, then the Kannada for a lake, then say \emph{guhe} again
\end{itemize}

\begin{tcolorbox}[breakable,colback=teal!4,colframe=teal!35!black,title={Before you move on}]
What does \kn{ಗುಹೆ} mean? (\textquotedblleft{}A cave\textquotedblright{}.) Say it once more.
\end{tcolorbox}
\section[\texorpdfstring{\kn{ಮೃಗಾಲಯ} (mṛgālaya)}{ಮೃಗಾಲಯ (mṛgālaya)}]{\kn{ಮೃಗಾಲಯ} (mṛgālaya) — a zoo}

\label{lesson:KA-C174-mrgalaya}



\begin{quote}
Before the new one: say the Kannada for a lake, then the Kannada for a cave.
\end{quote}

\subsection*{You'll want to know: \kn{ಮೃಗಾಲಯ}}
\textbf{\kn{ಮೃಗಾಲಯ}} — \emph{mṛgālaya} — \textquotedblleft{}a zoo\textquotedblright{}.

\begin{cousinweb}
\textbf{\kn{ಮೃಗಾಲಯ}} (\emph{mṛgālaya}) is Sanskrit \emph{mṛga}, \textquotedblleft{}wild animal, deer\textquotedblright{} (CDIAL 10264), and \emph{ālaya}, \textquotedblleft{}house\textquotedblright{} (CDIAL 1366): \textquotedblleft{}a house of animals\textquotedblright{}.
\end{cousinweb}

\begin{grammarlens}[title={What you've built}]
One more word for this chapter.
\end{grammarlens}

\subsection*{Your turn}
\begin{itemize}
\raggedright
  \item \emph{Say it:} \emph{mṛgālaya}
  \item \emph{Say it:} \emph{mṛgālaya}, once more
  \item \emph{Say it:} say \emph{mṛgālaya}
  \item \emph{Recall:} say the Kannada for a lake, then the Kannada for a cave, then say \emph{mṛgālaya} again
\end{itemize}

\begin{tcolorbox}[breakable,colback=teal!4,colframe=teal!35!black,title={Before you move on}]
What does \kn{ಮೃಗಾಲಯ} mean? (\textquotedblleft{}A zoo\textquotedblright{}.) Say it once more.
\end{tcolorbox}
\section[\texorpdfstring{\kn{ಅರಮನೆ} (aramane)}{ಅರಮನೆ (aramane)}]{\kn{ಅರಮನೆ} (aramane) — a palace}

\label{lesson:KA-C174-aramane}



\begin{quote}
Before the new one: say the Kannada for a cave, then the Kannada for a zoo.
\end{quote}

\subsection*{You'll want to know: \kn{ಅರಮನೆ}}
\textbf{\kn{ಅರಮನೆ}} — \emph{aramane} — \textquotedblleft{}a palace\textquotedblright{}.

\begin{cousinweb}
\textbf{\kn{ಅರಮನೆ}} (\emph{aramane}) is \textquotedblleft{}the king's house\textquotedblright{}. DEDR 201 lists \emph{ara-mane}, \textquotedblleft{}royal dwelling, palace\textquotedblright{}, itself. Its first part is \emph{arasu}, \textquotedblleft{}king\textquotedblright{}, which DEDR derives from Sanskrit \emph{rājan}, \textquotedblleft{}king\textquotedblright{}; the second is native \emph{mane}, \textquotedblleft{}house\textquotedblright{} (DEDR 4776).
\end{cousinweb}

\begin{grammarlens}[title={What you've built}]
One more word for this chapter.
\end{grammarlens}

\subsection*{Your turn}
\begin{itemize}
\raggedright
  \item \emph{Say it:} \emph{aramane}
  \item \emph{Say it:} \emph{aramane}, once more
  \item \emph{Say it:} say \emph{aramane}
  \item \emph{Recall:} say the Kannada for a cave, then the Kannada for a zoo, then say \emph{aramane} again
\end{itemize}

\begin{tcolorbox}[breakable,colback=teal!4,colframe=teal!35!black,title={Before you move on}]
What does \kn{ಅರಮನೆ} mean? (\textquotedblleft{}A palace\textquotedblright{}.) Say it once more.
\end{tcolorbox}
